% Independent preview; original geometry and numeric relationships retained.
% Same five-node binary star under two elimination orders.
\begin{tikzpicture}[
  x=1mm,y=1mm,font=\figurefont\footnotesize,
  vnode/.style={circle,draw=FigureInk,fill=FigurePaper,
    minimum size=7.5mm,inner sep=1pt},
  removed/.style={vnode,draw=FigureRed,fill=FigureRedFill,
    text=FigureInk},
  edge/.style={draw=FigureGrid,line width=.65pt},
  fill edge/.style={draw=FigureRed,line width=1pt},
  step arrow/.style={-{Stealth[length=1.8mm]},draw=LancetGreen,line width=.8pt}
]
  \node[font=\figurefont\small\mdseries] at (38,7) {（a）好顺序：先消去叶节点};
  \node[font=\figurefont\small\mdseries] at (123,7) {（b）坏顺序：先消去中心};

  \begin{scope}
    \node[vnode] (c1) at (18,-23) {$C$};
    \node[removed] (a1) at (18,-8) {$A$};
    \node[vnode] (b1) at (3,-23) {$B$};
    \node[vnode] (d1) at (33,-23) {$D$};
    \node[vnode] (e1) at (18,-38) {$E$};
    \foreach \v in {a1,b1,d1,e1}{\draw[edge] (c1)--(\v);}
    \node[anchor=north,align=center,text=FigureMuted,font=\FigureNoteFont\footnotesize] at (18,-44)
      {先消去$A$\\桶作用域$\{A,C\}$};

    \draw[step arrow] (38,-23)--(49,-23);

    \node[vnode] (c2) at (61,-23) {$C$};
    \node[vnode] (b2) at (49,-23) {$B$};
    \node[vnode] (d2) at (73,-23) {$D$};
    \node[vnode] (e2) at (61,-36) {$E$};
    \foreach \v in {b2,d2,e2}{\draw[edge] (c2)--(\v);}
    \node[anchor=north,align=center,text=LancetGreen] at (61,-44)
      {不产生填充边\\最大桶：$2$个变量};
    \node[anchor=north,align=center] at (38,-58)
      {$A\to B\to D\to E\to C$\quad
       二元表最多$2^2=4$项};
  \end{scope}

  \draw[FigureGrid,line width=.5pt] (80,5)--(80,-64);

  \begin{scope}
    \node[removed] (c3) at (102,-23) {$C$};
    \node[vnode] (a3) at (102,-8) {$A$};
    \node[vnode] (b3) at (87,-23) {$B$};
    \node[vnode] (d3) at (117,-23) {$D$};
    \node[vnode] (e3) at (102,-38) {$E$};
    \foreach \v in {a3,b3,d3,e3}{\draw[edge] (c3)--(\v);}
    \node[anchor=north,align=center,text=FigureMuted,font=\FigureNoteFont\footnotesize] at (102,-44)
      {先消去$C$\\桶含全部$5$个变量};

    \draw[step arrow] (122,-23)--(133,-23);

    \node[vnode] (a4) at (145,-8) {$A$};
    \node[vnode] (b4) at (133,-23) {$B$};
    \node[vnode] (d4) at (157,-23) {$D$};
    \node[vnode] (e4) at (145,-38) {$E$};
    \foreach \u/\v in {a4/b4,a4/d4,a4/e4,b4/d4,b4/e4,d4/e4}{
      \draw[fill edge] (\u)--(\v);
    }
    \node[anchor=north,align=center,text=FigureInk] at (145,-44)
      {$4$个邻居补成团\\产生$6$条填充边};
    \node[anchor=north,align=center] at (123,-58)
      {$C\to A\to B\to D\to E$\quad
       首桶访问$2^5=32$项};
  \end{scope}
\end{tikzpicture}
